\scanpage{96}
On va définir un foncteur fondamental \inserted{dit \uline{de Néron}}
\[
\mathcal N:C(S)\longrightarrow C'(k).
\]
Pour ceci, si $F$ est un faisceau de Weil sur $S$, on définit un faisceau
de Weil sur $\mathcal T_k$ en prenant, pour tt anneau \unsure{comm.} $A$ sur $k$
\inserted{\[
\simeq F(\varinjlim_{P}\mathbb V(A^P))
=F(\varprojlim_{P}\operatorname{Spec}(\mathbb V(A^P)))
\qquad\inserted{\operatorname{Spec}(\varinjlim_{P}\mathbb V(A^P))}
\]}
\[
\begin{aligned}
\mathcal N(F)(A)&=\varinjlim_{P}
 F\left(\unsure{\operatorname{Spec}}\left(\mathbb V
 \left(\underbrace{\prod_{i\in I_P} A_i^P}_{A^P}\right)\right)\right)\\
 &=\underline{\lim}_{P}\prod_{i\in I_P}F(\mathbb V(A_i^P)).
\end{aligned}
\]
\margin{N.B. On a en fait une ``application continue''
$\mathcal T'_k\longrightarrow\mathcal T_S$, \unsure{correspondant}
aux \unsure{pro-objets}, d'où le foncteur image directe de faisceaux, et
ses propriétés ``formelles''\ldots ! ? !}
\figure{\begin{tikzpicture}[x=1cm,y=1cm]
\draw(0,5)..controls(1.4,6.5)and(1.6,5)..(2.2,4.5)..controls(3,3.3)and(1,1)..(-.2,0);
\end{tikzpicture}}

$P$ parcourt les ``partitions constructibles finies'' de $A$, et où pour
une telle $P$, $I_P$ est son ensemble d'indices, et $A_i^P$ les composantes
correspondantes de $A$. En d'autres termes, un \unsure{quasi-schéma}
\unsure{flasque} \inserted{cpt} étant considéré comme un pro-(schéma affine)
sur $k$, on le transforme par $\mathbb V$ en un pro (schéma affine) sur $S$,
et on lui applique le foncteur $F$.
\deleted{[N.B. $\mathcal N(F)$ est compatible \unsure{avec} les limites
inductives d'anneaux \unsure{donc si on y tient}, on peut se limiter}
\margin{Non, car le foncteur $\mathbb V:k\text{-alg}\longrightarrow V\text{-alg.}$
ne commute pas aux limites inductives\ldots}
\mspage{95}

